\documentclass{article}
\usepackage{amsmath, amssymb, amsthm}
\usepackage{hyperref}
\usepackage{geometry}
\geometry{margin=1in}

\title{Erd\H{o}s Problem \#101}
\date{Last updated: 2025-08-31}
\author{Source: erdosproblems.com}

\begin{document}
\maketitle

\noindent\textbf{Status:} OPEN \quad \textbf{Prize: \$100}

\noindent\textbf{Tags:} geometry

\medskip
\noindent\textbf{Problem Statement:}

Given $n$ points in $\mathbb{R}^2$, no five of which are on a line, the number of lines containing four points is $o(n^2)$.

\bigskip
\noindent\textbf{Remarks:}

There are examples of sets of $n$ points with $\sim n^2/6$ many collinear triples and no four points on a line. Such constructions are given by Burr, Gr\"{u}nbaum, and Sloane \cite{BGS74} and F\"{u}redi and Pal\'{a}sti \cite{FuPa84}. 

Gr\"{u}nbaum \cite{Gr76} constructed an example with $\gg n^{3/2}$ such lines. Erd\H{o}s speculated this may be the correct order of magnitude. This is false: Solymosi and Stojakovi\'{c} \cite{SoSt13} have constructed a set with no five on a line and at least\[n^{2-O(1/\sqrt{\log n})}\]many lines containing exactly four points.

See also [102] and [669]. A generalisation of this problem is asked in [588].

This problem is Problem 71 on Green's open problems list.

\bigskip
\noindent\textbf{References:}

[BGS74] Burr, Stefan A. and Gr\"{u}nbaum, Branko and Sloane, N. J. A., The orchard problem. Geometriae Dedicata (1974), 397-424.
 
 [FuPa84] F\"{u}redi, Z. and Pal\'{a}sti, I., Arrangements of lines with a large number of triangles. Proc. Amer. Math. Soc. (1984), 561-566.
 
 [Gr76] Gr\"{u}nbaum, Branko, New views on some old questions of combinatorial geometry. Colloquio Internazionale sulle Teorie Combinatorie
(Roma, 1973), Tomo I (1976), 451-468.
 
 [SoSt13] Solymosi, J\'{o}zsef and Stojakovi\'C, Milo\vS, Many collinear {$k$}-tuples with no {$k+1$} collinear points. Discrete Comput. Geom. (2013), 811-820.

\bigskip
\noindent\small{Source: \url{https://www.erdosproblems.com/101}}
\end{document}
